	\subsection*{Question 1}
		$\np{45310}=\np{4,5310}\times \np{10000}=\np{4,531}\times 10^4$.

		\textbf{Réponse B}.

	\subsection*{Question 2}
		Dans l'expression \quad $(4x – 3)(4x + 3)$,

		on reconnaît le premier membre de l'identité remarquable \quad $(a+b)(a-b)=a^2-b^2$.

		En identifiant $a=4x$ et $b=3$, on obtient : \quad $(4x)^2-3^2=16x^2-9$.

		\textbf{Réponse C}.

	\subsection*{Question 3}
		Le volume d'un pavé droit est le produit de sa longueur par sa largeur par sa hauteur.

		On a donc :\quad $\mathcal{V}=4,5\times 4\times 10=\np[cm^3]{180}$.

		\textbf{Réponse A}.

	\subsection*{Question 4}
		Le critère de divisibilité par 9, c'est de voir si la somme des chiffres du nombre est divisible par 9.

		$2+0+2+5=9 = 9 \times 1$ \quad donc $\np{2025}$ est divisible par 9.

		$2+0+2+6=10 $\quad donc \quad $\np{2026}$ n'est pas divisible par 9.

		Donc \textbf{Réponse B}.

	\subsection*{Question 5}
		Effectuons la conversion :
		\quad
		$\begin{aligned}[t]
		\np[km]{9}  &\to  \np[min]{45} &\text{ en divisant par 3}\\
		\np[km]{3}  &\to  \np[min]{15} & \text{ en multipliant par 4}\\
		\np[km]{12} &\to  \np[min]{60} & \\
		\end{aligned}$

		Et \np[min]{60}, c'est une heure, donc sa vitesse est de 12 km/h.

	\subsection*{Question 6}
		Comme les secteurs sont de taille égale, on est en situation d'équiprobabilité, la probabilité est donc obtenue en divisant le nombre d'issues favorables par le nombre d'issues total.

		Il y a 2 casques audio sur 10 secteurs, donc la probabilité est de $\dfrac{2}{10}$, c'est-à-dire, après simplification $\dfrac{1}{5}$.

	\subsection*{Question 7}

		\emph{Méthode 1 :}

		$60\times \dfrac{10}{100}=6$, \quad 10\,\% de 60\,\euro{} représentent donc 6\,\euro{}.

		$60-6=54$, \quad son nouveau prix après une baisse de 10\,\% est de 54 \euro{}.

		\emph{Méthode 2 :}

		Baisser de 10\,\%, c'est multiplier par : \quad $c = 1 - \dfrac{10}{100} = 0,9$.

		$60c = 60 \times 0,9 = 6 \times 9 = 54$,\quad son nouveau prix après une baisse de 10\,\% est de 54 \euro{}.

	\subsection*{Question 8}
		La somme des angles dans un triangle vaut 180\textdegree.

		$180 = 90+ 40 + x$\quad soit \quad $x =180 -(90+40)=50$.

		L'angle $\widehat{\mathrm{BAC}}$ a donc une mesure de 50\textdegree.

	\subsection*{Question 9}

		\begin{enumerate}[label=\textbf{\alph*.},left=5mm]
		\item Le nombre d'élèves ayant participé à ce contrôle est la somme des effectifs des notes :

		$3+4+0+4+5+6+0+0+3+0+2+1=28$,\quad il y a 28 élèves.

		\item $28 \div 2 = 14$,\quad la moitié de l'effectif est 14.

		L'effectif total est pair, donc la médiane est la moyenne des deux valeurs centrales : la 14\up{e} et la 15\up{e} dans la série ordonnée.

		Dans l'ordre croissant :
		\begin{itemize}
			\item Les trois premières valeurs sont des 7;
			\item les quatre valeurs suivantes, de la 4\up{e} à la 7\up{e}, sont des 8;
			\item les quatre valeurs suivantes, de la 8\up{e} à la 11\up{e} sont des 10;
			\item les cinq valeurs suivantes, de la 12\up{e} à la 16\up{e} sont des 11.
		\end{itemize}

		Ainsi, la 14\up{e} et la valeur dans l'ordre croissant sont égales à 11.

		$\dfrac{11+11}{2}=11$,\quad a médiane est donc 11.
		\end{enumerate}
